\ifdefined\gkver\else\def\gkver{student}\fi
\documentclass[\gkver]{gaokaozhenti}
\begin{document}

\chapter{2010年全国理综Ⅰ}
%% paperType: 理综物理部分

\begin{enumerate}
\item
%% number: 14
%% paperName: 2010年全国理综Ⅰ第 14 题 6 分
%% typeId: 1
%% type: 单选题
%% chapter: 原子和原子核
%% point: 天然放射性
%% method: 无
%% score: 6
%% degree: 650
%% duplicateId: 0
%% body:
原子核 \ce{^{238}_{92}U} 经放射性衰变①变为原子核 \ce{^{234}_{90}Th}，继而经放射性衰变②变为原子核 \ce{^{234}_{91}Pa}，再经放射性衰变③变为原子核 \ce{^{234}_{92}U}。放射性衰变①、②和③依次为\xzanswer{A}
\fourchoices[answer=A]
{$\alpha$ 衰变、$\beta$ 衰变和 $\beta$ 衰变}
{$\beta$ 衰变、$\alpha$ 衰变和 $\beta$ 衰变}
{$\beta$ 衰变、$\beta$ 衰变和 $\alpha$ 衰变}
{$\alpha$ 衰变、$\beta$ 衰变和 $\alpha$ 衰变}
%% answer: A
%% memo:
\memoanswer{
$\ce{^{238}_{92}U -> ^{234}_{90}Th}$，质量数少 4，电荷数少 2，说明衰变①为 $\alpha$ 衰变：$\ce{^{238}_{92}U -> ^{234}_{90}Th + ^{4}_{2}He}$。$\ce{^{234}_{90}Th -> ^{234}_{91}Pa}$，质子数加 1，说明衰变②为 $\beta$ 衰变，中子转化成质子：$\ce{^{234}_{90}Th -> ^{234}_{91}Pa + ^{0}_{-1}e}$。$\ce{^{234}_{91}Pa -> ^{234}_{92}U}$，质子数加 1，说明衰变③为 $\beta$ 衰变，中子转化成质子：$\ce{^{234}_{91}Pa -> ^{234}_{92}U + ^{0}_{-1}e}$。
}

\item
%% number: 15
%% paperName: 2010年全国理综Ⅰ第 15 题 6 分
%% typeId: 1
%% type: 单选题
%% chapter: 运动和力
%% point: 瞬时加速度
%% method: 无
%% score: 6
%% degree: 450
%% duplicateId: 0
%% body:
如右图，轻弹簧上端与一质量为 $m$ 的木块 1 相连，下端与另一质量为 $M$ 的木块 2 相连，整个系统置于水平放置的光滑木板上，并处于静止状态。现将木板沿水平方向突然抽出，设抽出后的瞬间，木块 1、2 的加速度大小分别为 $a_{1}$、$a_{2}$，重力加速度大小为 $g$。则有\xzanswer{C}
\onepicture[w=3.5cm]{figs/15.png}
\fourchoices[answer=C]
{$a_{1} = 0$，$a_{2} = g$}
{$a_{1} = g$，$a_{2} = g$}
{$a_{1} = 0$，$a_{2} = \frac{M + m}{M}g$}
{$a_{1} = g$，$a_{2} = \frac{M + m}{M}g$}
%% answer: C
%% memo:
\memoanswer{
在抽出木板的瞬间，弹簧对 1 的支持力和对 2 的压力并未改变。对 1 物体受重力和支持力，$mg = F$，$a_{1} = 0$。对 2 物体受重力和压力，根据牛顿第二定律 $a_{2} = \frac{F + Mg}{M} = \frac{M + m}{M}g$。
}

\item
%% number: 16
%% paperName: 2010年全国理综Ⅰ第 16 题 6 分
%% typeId: 1
%% type: 单选题
%% chapter: 电场
%% point: 电势能电势差电场力做功
%% method: 无
%% score: 6
%% degree: 650
%% duplicateId: 0
%% body:
关于静电场，下列结论普遍成立的是\xzanswer{C}
\fourchoices[answer=C]
{电场强度大的地方电势高，电场强度小的地方电势低}
{电场中任意两点之间的电势差只与这两点的场强有关}
{在正电荷或负电荷产生的静电场中，场强方向都指向电势降低最快的方向}
{将正点电荷从场强为零的一点移动到场强为零的另一点，电场力做功这零}
%% answer: C
%% memo:
\memoanswer{
在正电荷的电场中，离正电荷近，电场强度大，电势高，离正电荷远，电场强度小，电势低；而在负电荷的电场中，离负电荷近，电场强度大，电势低，离负电荷远，电场强度小，电势高，A 错误。电势差的大小决定于两点间距和电场强度，B 错误；沿电场方向电势降低，而且速度最快，C 正确；场强为零，电势不一定为零，如从带正电荷的导体球上将正电荷移动到另一带负电荷的导体球上，电场力做正功，D 错误。
}

\item
%% number: 17
%% paperName: 2010年全国理综Ⅰ第 17 题 6 分
%% typeId: 2
%% type: 多选题
%% chapter: 电磁感应
%% point: 切割磁感线电动势
%% method: 无
%% score: 6
%% degree: 650
%% duplicateId: 0
%% body:
某地的地磁场磁感应强度的竖直分量方向向下，大小为 $4.5\times10^{-5}\UT$。一灵敏电压表连接在当地入海河段的两岸，河宽 $100\Um$，该河段涨潮和落潮时有海水（视为导体）流过。设落潮时，海水自西向东流，流速为 $2\Ums$。下列说法正确的是\xzanswer{BD}
\fourchoices[answer=BD]
{电压表记录的电压为 $5\UmV$}
{电压表记录的电压为 $9\UmV$}
{河南岸的电势较高}
{河北岸的电势较高}
%% answer: BD
%% memo:
\memoanswer{
海水在落潮时自西向东流，该过程可以理解为：自西向东运动的导体棒在切割竖直向下的磁场。根据右手定则，右岸即北岸是正极、电势高，南岸电势低，D 对 C 错。根据法拉第电磁感应定律 $E = BLv = 4.5\times10^{-5}\times100\times2 = 9\times10^{-3}\UV$，B 对 A 错。
}

\item
%% number: 18
%% paperName: 2010年全国理综Ⅰ第 18 题 6 分
%% typeId: 1
%% type: 单选题
%% chapter: 曲线运动
%% point: 平抛运动
%% method: 无
%% score: 6
%% degree: 450
%% duplicateId: 0
%% body:
一水平抛出的小球落到一倾角为 $\theta$ 的斜面上时，其速度方向与斜面垂直，运动轨迹如右图中虚线所示。小球在竖直方向下落的距离与在水平方向通过的距离之比为\xzanswer{D}
\onepicture[w=4.5cm]{figs/18.png}
\fourchoices[answer=D]
{$\tan\theta$}
{$2\tan\theta$}
{$\frac{1}{\tan\theta}$}
{$\frac{1}{2\tan\theta}$}
%% answer: D
%% memo:
\memoanswer{
如图平抛的末速度与竖直方向的夹角等于斜面倾角 $\theta$，有：

$\tan\theta = \frac{v_{0}}{gt}$

则下落高度与水平射程之比为

$\frac{y}{x} = \frac{\frac{1}{2}gt^{2}}{v_{0}t} = \frac{gt}{2v_{0}} = \frac{1}{2\tan\theta}$

\onepicture[w=4.5cm]{figs/18_an.png}

D 正确。
}

\item
%% number: 19
%% paperName: 2010年全国理综Ⅰ第 19 题 6 分
%% typeId: 2
%% type: 多选题
%% chapter: 气体性质
%% point: 分子力
%% method: 无
%% score: 6
%% degree: 450
%% duplicateId: 0
%% body:
右图为两分子系统的势能 $E_{p}$ 与两分子间距离 $r$ 的关系曲线。下列说法正确的是\xzanswer{BC}
\onepicture[w=5cm]{figs/19.png}
\fourchoices[answer=BC]
{当 $r$ 大于 $r_{1}$ 时，分子间的作用力表现为引力}
{当 $r$ 小于 $r_{1}$ 时，分子间的作用力表现为斥力}
{当 $r$ 等于 $r_{2}$ 时，分子间的作用力为零}
{当 $r$ 由 $r_{1}$ 变到 $r_{2}$ 的过程中，分子间的作用力做负功}
%% answer: BC
%% memo:
\memoanswer{
分子间距等于 $r_{0}$ 时分子势能最小，即 $r_{0} = r_{2}$。当 $r$ 小于 $r_{1}$ 时分子力表现为斥力；当 $r$ 大于 $r_{1}$ 小于 $r_{2}$ 时分子力表现为斥力；当 $r$ 大于 $r_{2}$ 时分子力表现为引力，A 错、B、C 对。在 $r$ 由 $r_{1}$ 变到 $r_{2}$ 的过程中，分子斥力做正功，分子势能减小，D 错误。
}

\item
%% number: 20
%% paperName: 2010年全国理综Ⅰ第 20 题 6 分
%% typeId: 1
%% type: 单选题
%% chapter: 光学
%% point: 光的折射
%% method: 无
%% score: 6
%% degree: 450
%% duplicateId: 0
%% body:
某人手持边长为 $6\Ucm$ 的正方形平面镜测量身后一棵树的高度。测量时保持镜面与地面垂直，镜子与眼睛的距离为 $0.4\Um$。在某位置时，他在镜中恰好能够看到整棵树的像；然后他向前走了 $6.0\Um$，发现用这个镜子长度的 $5/6$ 就能看到整棵树的像。这棵树的高度约为\xzanswer{B}
\fourchoices[answer=B]
{$4.0\Um$}
{$4.5\Um$}
{$5.0\Um$}
{$5.5\Um$}
%% answer: B
%% memo:
\memoanswer{
如图是恰好看到树时的反射光路，设树高为 $H$，树到镜的距离为 $L$，由图中的三角形可得

$\frac{H}{0.06\Um} = \frac{L + 0.4\Um}{0.4\Um}$

人离树越远，视野越大，看到树所需镜面越小，同理有

$\frac{H}{0.05\Um} = \frac{L + 0.4\Um + 6\Um}{0.4\Um}$

以上两式解得 $L = 29.6\Um$，$H = 4.5\Um$。

\onepicture[w=6cm]{figs/20_an.png}
}

\item
%% number: 21
%% paperName: 2010年全国理综Ⅰ第 21 题 6 分
%% typeId: 1
%% type: 单选题
%% chapter: 机械振动
%% point: 弹簧振子
%% method: 无
%% score: 6
%% degree: 450
%% duplicateId: 0
%% body:
一简谐振子沿 $x$ 轴振动，平衡位置在坐标原点。$t = 0$ 时刻振子的位移 $x = -0.1\Um$；$t = \frac{4}{3}\Us$ 时刻 $x = 0.1\Um$；$t = 4\Us$ 时刻 $x = 0.1\Um$。该振子的振幅和周期可能为\xzanswer{A}
\fourchoices[answer=A]
{$0.1\Um$，$\frac{8}{3}\Us$}
{$0.1\Um$，$8\Us$}
{$0.2\Um$，$\frac{8}{3}\Us$}
{$0.2\Um$，$8\Us$}
%% answer: A
%% memo:
\memoanswer{
在 $t = \frac{4}{3}\Us$ 和 $t = 4\Us$ 两时刻振子的位移相同。第一种情况是此时间差是周期的整数倍：$4 - \frac{4}{3} = nT$，当 $n = 1$ 时 $T = \frac{8}{3}\Us$。在 $\frac{4}{3}\Us$ 的半个周期内振子的位移由负的最大变为正的最大，所以振幅是 $0.1\Um$，A 正确。

第二种情况是此时间差不是周期的整数倍，则 $\left(\frac{4}{3} - 0\right) + \left(4 - \frac{4}{3}\right) = nT + \frac{T}{2}$，当 $n = 0$ 时 $T = 8\Us$，且由于 $\Delta t_{2}$ 是 $\Delta t_{1}$ 的二倍，说明振幅是该位移的二倍为 $0.2\Um$，对应选项 D。

\onepicture[w=3cm]{figs/21_an.png}
}

\item
%% number: 22
%% paperName: 2010年全国理综Ⅰ第 22 题 6 分
%% typeId: 4
%% type: 实验题
%% chapter: 曲线运动
%% point: 圆周运动的应用
%% method: 无
%% score: 6
%% degree: 450
%% duplicateId: 0
%% body:
图\ref{2010全国理综Ⅰ22a}是利用激光测转速的原理示意图，图中圆盘可绕固定轴转动，盘边缘侧面上有一小段涂有很薄的反光材料。当盘转到某一位置时，接收器可以接受到反光涂层所反射的激光束，并将所收到的光信号转变成电信号，在示波器显示屏上显示出来（如图\ref{2010全国理综Ⅰ22b}所示）。

\twopicture[numA=1, numB=2, labelA=2010全国理综Ⅰ22a, labelB=2010全国理综Ⅰ22b, wA=7cm, wB=7cm]
{figs/22a.png}{figs/22b.png}

\begin{enumerate}
	\item
	若图\ref{2010全国理综Ⅰ22b}中示波器显示屏横向的每大格（5 小格）对应的时间为 $5.00\times10^{-2}\Us$，则圆盘的转速为\tkanswer{4.55} $\Urs$。（保留 3 位有效数字）

	\item
	若测得圆盘直径为 $10.20\Ucm$，则可求得圆盘侧面反光涂层的长度为\tkanswer{1.46} $\Ucm$。（保留 3 位有效数字）

\end{enumerate}
%% answer: 
\jdanswer{
\begin{enumerate}
	\item
	$4.55\Urs$

	\item
	$1.46\Ucm$

\end{enumerate}
}
%% memo:
\memoanswer{
（1）从图 2 可知圆盘转一圈的时间在横坐标上显示 22 格，由题意知图 2 中横坐标上每格表示 $1.00\times10^{-2}\Us$，所以圆盘转动的周期是 $0.22\Us$，则转速为 $4.55\Urs$。

（2）反光引起的电流图像在图 2 中横坐标上每次一格，说明反光涂层的长度占圆盘周长的 22 分之一，为 $\frac{2\pi r}{22} = \frac{2\times3.14\times5.10}{22}\Ucm = 1.46\Ucm$。
}

\item
%% number: 23
%% paperName: 2010年全国理综Ⅰ第 23 题 12 分
%% typeId: 4
%% type: 实验题
%% chapter: 电路
%% point: 电路实验
%% method: 无
%% score: 12
%% degree: 450
%% duplicateId: 0
%% body:
一电流表的量程标定不准确，某同学利用图 1 所示电路测量该电流表的实际量程 $I_{m}$。所用器材有：

量程不准的电流表 A$_{1}$，内阻 $r_{1} = 10.0\UO$，量程标称为 $5.0\UmA$；

标准电流表 A$_{2}$，内阻 $r_{2} = 45.0\UO$，量程标称为 $1.0\UmA$；

标准电阻 $R_{1}$，阻值 $10.0\UO$；

滑动变阻器 $R$，总电阻约为 $300.0\UO$；

电源 $E$，电动势 $3.0\UV$，内阻不计；

保护电阻 $R_{2}$；开关 S；导线。

回答下列问题：

\begin{enumerate}
	\item
	在答题纸上（图 2 所示）的实物图上画出连线。

	\drawpicanswer{%
	\onepicture[w=9cm, align=right]{figs/23b.png}%
	}{%
	\onepicture[w=9cm, align=right]{figs/23_an.png}%
	}

	\item
	开关 S 闭合前，滑动变阻器的滑动端 c 应滑至\tkanswer{b}端。

	\item
	开关 S 闭合后，调节滑动变阻器的滑动端，使电流表 A$_{1}$ 满偏，若此时电流表 A$_{2}$ 的读数为 $I_{2}$，则 A$_{1}$ 的量程 $I_{m} =$\tkanswer{$\frac{R_{1} + r_{2}}{r_{1}}I_{2} = 5.5I_{2}$}。

	\item
	若测量时，A$_{1}$ 未调到满偏，两电流表的示数如图 3 所示，从图中读出 A$_{1}$ 的示数 $I_{1} =$\tkanswer{$3.00\UmA$}，A$_{2}$ 的示数 $I_{2} =$\tkanswer{$0.660\UmA$}；由读出的数据计算得 $I_{m} =$\tkanswer{$6.05\UmA$}。（保留 3 位有效数字）

	\item
	写出一条提高测量准确度的建议：\tkanswer{多次测量求平均}。

\end{enumerate}

\onepicture[align=right, w=6.5cm]{figs/23a.png}
\onepicture[w=12cm]{figs/23c.png}
%% answer: 
\jdanswer{
\begin{enumerate}
	\item
	连线见答案图。

	\item
	b

	\item
	$I_{m} = \frac{R_{1} + r_{2}}{r_{1}}I_{2} = 5.5I_{2}$

	\item
	$3.00\UmA$，$0.660\UmA$，$6.05\UmA$

	\item
	多次测量求平均（其它合理答案同样给分，如：测量时，电流表指针偏转大于满刻度的 1/2）

\end{enumerate}
}
%% memo:
\memoanswer{
（1）连线见答案图。

（2）在滑动变阻器的限流接法中，在接通开关前需要将滑动触头滑动到阻值最大端。

（3）闭合开关调节滑动变阻器使待测表满偏，流过的电流为 $I_{m}$。根据并联电路电压相等有 $I_{m}r_{1} = I_{2}(r_{2} + R_{1})$，得 $I_{m} = \frac{I_{2}(R_{1} + r_{2})}{r_{1}}$。

（4）待测表未满偏有 $I_{1}r_{1} = I_{2}(r_{2} + R_{1})$，将 A$_{2}$ 的示数 $0.66\UmA$ 和其他已知条件代入有

$I_{1} = \frac{I_{2}(R_{1} + r_{2})}{r_{1}} = \frac{0.66\times(10 + 45)}{10}\UmA = 3.63\UmA$

但图中 A$_{1}$ 的示数为 $3.0\UmA$、量程为 $5.0\UmA$，根据电流表的刻度是均匀的，则准确量程为 $6.05\UmA$。
}

\item
%% number: 24
%% paperName: 2010年全国理综Ⅰ第 24 题 15 分
%% typeId: 6
%% type: 计算题
%% chapter: 直线运动
%% point: 运动图象
%% method: 无
%% score: 15
%% degree: 650
%% duplicateId: 0
%% body:
汽车由静止开始在平直的公路上行驶，$0\sim60\Us$ 内汽车的加速度随时间变化的图线如右图所示。

\begin{enumerate}
	\item
	画出汽车在 $0\sim60\Us$ 内的 $v-t$ 图线；

	\item
	求这 $60\Us$ 内汽车行驶的路程。

\end{enumerate}
\onepicture[align=right, w=6cm]{figs/24.png}
%% answer: 
\jdanswer{
\begin{enumerate}
	\item
	速度图像如图所示。

	\item
	$900\Um$

\end{enumerate}
}
%% memo:
\memoanswer{
（1）由加速度图像可知前 $10\Us$ 汽车匀加速运动，后 $20\Us$ 汽车匀减速运动恰好停止，因为图像的面积表示速度的变化，此两段的面积相等。最大速度为 $20\Ums$。所以速度图像为下图。

\onepicture[w=6cm]{figs/24_an.png}

然后利用速度图像的面积求出位移。

（2）汽车运动的面积为匀加速、匀速、匀减速三段的位移之和。

$s = s_{1} + s_{2} + s_{3} = 10\times10 + 30\times20 + 10\times20 = 900\Um$
}

\item
%% number: 25
%% paperName: 2010年全国理综Ⅰ第 25 题 18 分
%% typeId: 6
%% type: 计算题
%% chapter: 曲线运动
%% point: 双星问题
%% method: 无
%% score: 18
%% degree: 450
%% duplicateId: 0
%% body:
如右图，质量分别为 $m$ 和 $M$ 的两个星球 A 和 B 在引力作用下都绕 O 点做匀速圆周运动，星球 A 和 B 两者中心之间的距离为 $L$。已知 A、B 的中心和 O 三点始终共线，A 和 B 分别在 O 的两侧。引力常数为 $G$。

\begin{enumerate}
	\item
	求两星球做圆周运动的周期；

	\item
	在地月系统中，若忽略其他星球的影响，可以将月球和地球看成上述星球 A 和 B，月球绕其轨道中心运行的周期为 $T_{1}$。但在近似处理问题时，常常认为月球是绕地心做圆周运动的，这样算得的运行周期记为 $T_{2}$。已知地球和月球的质量分别为 $5.98\times10^{24}\Ukg$ 和 $7.35\times10^{22}\Ukg$。求 $T_{2}$ 与 $T_{1}$ 两者平方之比。（结果保留 3 位小数）

\end{enumerate}
\onepicture[align=right, w=4.5cm]{figs/25.png}
%% answer: 
\jdanswer{
\begin{enumerate}
	\item
	$T = 2\pi\sqrt{\frac{L^{3}}{G(M + m)}}$

	\item
	$1.012$

\end{enumerate}
}
%% memo:
\memoanswer{
（1）设 A、B 星做圆周运动的轨道半径分别为 $r$ 和 $R$，相互作用力为 $f$，运动周期为 $T$。它们之间的万有引力提供向心力，则 A 和 B 的向心力相等。且 A、B 和 O 三点始终共线，说明 A 和 B 有相同的角速度和周期。

则有 $m\omega^{2}r = M\omega^{2}R$，$r + R = L$，解得 $R = \frac{m}{m + M}L$，$r = \frac{M}{m + M}L$。

对 A 根据牛顿第二定律和万有引力定律得

$\frac{GMm}{L^{2}} = m\left(\frac{2\pi}{T}\right)^{2}\frac{M}{M + m}L$ ①

化简得

$T = 2\pi\sqrt{\frac{L^{3}}{G(M + m)}}$ ②

（2）将地月看成双星，由①得

$T_{1} = 2\pi\sqrt{\frac{L^{3}}{G(M + m)}}$ ③

将月球看作绕地心做圆周运动，根据牛顿第二定律和万有引力定律得

$\frac{GMm}{L^{2}} = m\left(\frac{2\pi}{T}\right)^{2}L$ ④

化简得

$T_{2} = 2\pi\sqrt{\frac{L^{3}}{GM}}$ ⑤

所以两种周期的平方比值为

$\left(\frac{T_{2}}{T_{1}}\right)^{2} = \frac{m + M}{M} = \frac{5.98\times10^{24} + 7.35\times10^{22}}{5.98\times10^{24}} = 1.012$
}

\item
%% number: 26
%% paperName: 2010年全国理综Ⅰ第 26 题 21 分
%% typeId: 6
%% type: 计算题
%% chapter: 磁场
%% point: 带电粒子在磁场中的运动
%% method: 无
%% score: 21
%% degree: 250
%% duplicateId: 0
%% body:
如图，在 $0\leq x\leq\sqrt{3}a$ 区域内存在与 $xy$ 平面垂直的匀强磁场，磁感应强度的大小为 $B$。在 $t = 0$ 时刻，一位于坐标原点的粒子源在 $xy$ 平面内发射出大量同种带电粒子，所有粒子的初速度大小相同，方向与 $y$ 轴正方向夹角分布在 $0\sim180^\circ$ 范围内。已知沿 $y$ 轴正方向发射的粒子在 $t = t_{0}$ 时刻刚好从磁场边界上 P（$\sqrt{3}a$，$a$）点离开磁场。求：

\begin{enumerate}
	\item
	粒子在磁场中做圆周运动的半径 $R$ 及粒子的比荷 $q/m$；

	\item
	此时刻仍在磁场中的粒子的初速度方向与 $y$ 轴正方向夹角的取值范围；

	\item
	从粒子发射到全部粒子离开磁场所用的时间。

\end{enumerate}
\onepicture[align=right, w=6cm]{figs/26.png}
%% answer: 
\jdanswer{
\begin{enumerate}
	\item
	$R = \frac{2\sqrt{3}}{3}a$，$\frac{q}{m} = \frac{2\pi}{3Bt_{0}}$

	\item
	速度与 $y$ 轴的正方向的夹角范围是 $60^\circ$ 到 $120^\circ$

	\item
	从粒子发射到全部离开所用时间为 $2t_{0}$

\end{enumerate}
}
%% memo:
\memoanswer{
（1）粒子沿 $y$ 轴的正方向进入磁场，从 P 点经过做 OP 的垂直平分线与 $x$ 轴的交点为圆心，根据直角三角形有

$R^{2} = a^{2} + (\sqrt{3}a - R)^{2}$ ①

解得

$R = \frac{2\sqrt{3}}{3}a$ ②

\onepicture[w=4.5cm]{figs/26_an1.png}

$\sin\theta = \frac{a}{R} = \frac{\sqrt{3}}{2}$，则粒子做圆周运动的圆心角为 $120^\circ$，周期为 $T = 3t_{0}$ ③

粒子做圆周运动的向心力由洛伦兹力提供，根据牛顿第二定律得

$Bqv = m\left(\frac{2\pi}{T}\right)^{2}R$，$v = \frac{2\pi R}{T}$

化简得

$\frac{q}{m} = \frac{2\pi}{3Bt_{0}}$ ④

（2）仍在磁场中的粒子其圆心角一定大于 $120^\circ$，这样粒子角度最小时从磁场右边界穿出；角度最大时从磁场左边界穿出。

角度最小时从磁场右边界穿出、圆心角 $120^\circ$，所经过圆弧的弦与（1）中相等穿出点如图，根据弦与半径、$x$ 轴的夹角都是 $30^\circ$，所以此时速度与 $y$ 轴的正方向的夹角是 $60^\circ$。

\onepicture[w=4cm]{figs/26_an2.png}

角度最大时从磁场左边界穿出，半径与 $y$ 轴的夹角是 $60^\circ$，则此时速度与 $y$ 轴的正方向的夹角是 $120^\circ$。

\onepicture[w=4cm]{figs/26_an3.png}

所以速度与 $y$ 轴的正方向的夹角范围是 $60^\circ$ 到 $120^\circ$。

（3）在磁场中运动时间最长的粒子的轨迹应该与磁场的右边界相切，在三角形中两个相等的腰为 $R = \frac{2\sqrt{3}}{3}a$，而它的高是

$h = \sqrt{3}a - \frac{2\sqrt{3}}{3}a = \frac{\sqrt{3}}{3}a$

半径与 $y$ 轴的夹角是 $30^\circ$，这种粒子的圆心角是 $240^\circ$，所用时间为 $2t_{0}$。

所以从粒子发射到全部离开所用时间为 $2t_{0}$。

\onepicture[w=3.5cm]{figs/26_an4.png}
}

\end{enumerate}

\end{document}
